\section{Bewertung des Testers}\label{sec:bewertung_des_testers}
%Positiv
Der Tester ist gut geeignet um diesen Algorithmus zu Testen. Er implementiert eine Logik zur Überprüfung der Genauigkeit und erkennt, welche Gipfel gefunden werden. Die damit zusammenhängende Umformulierung des Gipfelsuchproblem zu einem Klassifizierungsproblem hat geholfen um schnell zu erkennen zu können, wie gut der Algorithmus Gipfel erkennt. Diese Methodik hat das manuele Vergelichen der Algorithmusergebnisse überflüssig gemacht. Die in der falschen Gipfel in der Realdaten-Datei \todo{besserer Name für die datei und konsistenz dieses namens} haben bei der Diagnose von Fehlern geholfen. Dadurch konnten nicht gefundene Gipfel schneller auf einer Karte entdeckt werden. 

% Negativ
Die Wahl der Testdaten war nicht ideal. Es hätten weniger Geodaten des Copernikusprogramm verwendet werden sollen. Diese sind ungenauer als die des baden-württembergischen oder bayerischen Vermessungsamtes. Um diese Daten nutzen zu können hätten andere Kartenausschnitte gewählt werden müssen. Es mussten auch viele Realdaten manuell überprüft und ergänzt werden. Das hätte auch mit der Wahl einer höheren minimalen Prominenz und Dominanz verbessert werden könne, da dadurch nach weniger unbekannten Gipfeln auf den Kartenausschnitten gesucht werden würde. Das wäre im Schwarzwald nicht zielführend gewesen und somit hätten die untersuchetn Kartenausschnitte anders gewählt werden müssen. 